Un fil conductor molt llarg segueix l'eix $z$ i transporta un corrent $I = 2,00\ \mathrm{A}$. Quan arriba a l'altura de $z = 0,00\ \mathrm{cm}$ canvia de direcció i traça una circumferència en el pla $xy$, de radi $R = 2,00\ \mathrm{cm}$ i centrada al punt $(0, R, 0)$, i després continua per l'eix $z$.

\figura[0.4]{fig1.png}

\begin{apartat}{1,25}
Calculeu el vector i el mòdul del camp magnètic total al centre de la circumferència, és a dir, al punt $(0, R, 0)$.
\end{apartat}

\begin{apartat}{1,25}
Podem girar a voluntat l'espira respecte a l'eix $y$. En quina direcció hem d'orientar-la per a obtenir el mòdul del camp magnètic mínim i el mòdul del camp magnètic màxim? Trobeu aquests dos valors del mòdul del camp magnètic i especifiqueu el pla on hi ha l'espira i el sentit de gir del corrent en cada cas.
\end{apartat}

\dades{%
  El mòdul del camp magnètic creat per un fil infinit per on circula un corrent $I$ a una distància $r$ del fil és $B = \dfrac{\mu_0 I}{2\pi r}$.\\[2pt]
  El mòdul del camp magnètic creat al centre d'una espira de radi $R$ per on circula un corrent $I$ és $B = \dfrac{\mu_0 I}{2R}$.\\[2pt]
  $\mu_0 = 4\pi\times 10^{-7}\ \mathrm{T\,m\,A^{-1}}$.}
